% Experiment log and working notes.
% Record of every experiment run, every negative result, and every direction
% tried and dropped.
% Style rules: never use the em-dash; no italic text; -ize/-ization spelling
% (customization, factorization, optimization), except -wise suffixes, revise,
% exercise, noise/denoising, and file paths, which keep their own spelling.
\section{Experiment log and working notes}
\label{app:log}

This appendix is the running record of the project. It lists what was
measured, on which data and hardware, and what was tried and abandoned.
Negative results are kept on purpose: several of them decide what the paper
may and may not claim. Every number is traceable to an artifact in the code
repository; the path is given in each subsection.

\subsection{Three axes of customization}
\label{app:log:axes}

A solver is varied along three axes, and each experiment below moves one of
them. Naming them separately is what makes the negative results readable: a
recipe that wins on one family and loses on another has not failed, it has
located an axis-two boundary.

\begin{enumerate}
\item The solving algorithm. The base operator, the wrappers composed onto it
 (relaxation, anchoring, fixed and adaptive restarts, safeguarded
 acceleration), the metric, and the certificate cadence. This axis acts
 mainly on $N$, the number of iterations needed to certify.
\item The input problem family. Real data at several scales, described by the
 structure of the instance and by its conditioning. This axis does not make a
 recipe faster; it decides which recipe is right, and which baseline is even
 able to run.
\item The hardware. Execution structure, storage precision, and kernel and
 iteration fusion. This axis acts on $t$, the cost of one iteration.
\end{enumerate}

The product $T = N \times t$ is what a user experiences, and the certificate
cadence is the only gene that appears in both factors, which is why it is
treated separately throughout.

\subsection{What was run, and what was found}
\label{app:log:scoreboard}

Every row is measured against solvers written by their authors, on downloaded
real data except where marked generated, with acceptance being our own
float64 certificate in the original coordinates. Detail and artifact paths are
in the subsections that follow.

\begin{center}
\small
\begin{tabular}{p{2.5cm}p{2.6cm}p{2.2cm}p{6.1cm}}
\toprule
Family & Data and scale & Baseline & Outcome \\
\midrule
TV denoising & six images, $n=128$ to $496$ & our own XLA loop & Fusion wins on the hardware axis: $1.36\times$ to $2.63\times$ end to end on the RTX 4090, $1.15\times$ to $1.73\times$ on the Blackwell. Admission is bitwise identical at $\rho=1$. \\
\addlinespace
Elastic-net logistic & epsilon (LIBSVM), $10^5\times2000$, ten penalty settings & cuML, same GPU & $2.28\times$ ($95\%$ interval $[1.73, 3.01]$) on ten of ten for a repeated solve, and $0.06\times$ for a first solve on an empty cache. The margin is monotone in the penalty weight and crosses over near $0.05\lambda_{\max}$. \\
\addlinespace
Convex QP, Maros-Meszaros & 137 published instances & OSQP, SCS, Clarabel & No win. We certify 55 to 63 instances against OSQP's 123 and Clarabel's 132, and are $4.5\times$ to $5.8\times$ slower where both solve. Every instance factorizes cheaply. \\
\addlinespace
Convex QP, generated lasso & $n$ from $25{,}000$ to $750{,}000$ variables & OSQP, SCS, Clarabel & $1.98\times$ over Clarabel at $25{,}000$ and $6.04\times$ at $75{,}000$. At $250{,}000$ we finish in $102.7$ s while Clarabel is killed by the operating system at $28.9$ GB and OSQP does not certify in $5400$ s. Our own cost per iteration also degrades with size. \\
\addlinespace
Convex QP, generated SVM & $n$ from $3300$ to $110{,}000$ & OSQP, SCS, Clarabel & The ADMM solvers are beaten from about $11{,}000$ variables, and Clarabel never is: its lead grows with size, because the coupling constraints give an interior-point method a system that factorizes well. Measured before the timing correction described above and not repeated. \\
\addlinespace
Optimal transport, DOTmark & $32\times32$ and $64\times64$ & POT exact network simplex & Lost by $5\times$ to $32\times$ and abandoned. The exact simplex is very strong at these sizes, and the entropic scheme never certified on the original problem. \\
\bottomrule
\end{tabular}
\end{center}

Read along axis two, the pattern is that the wins are family dependent rather
than merely size dependent. A large instance is not sufficient: the generated
SVM family is large and is never won against Clarabel, while the generated
lasso family of the same size is, because the two give an interior-point
method very different systems to factorize. The one clean statement at the
largest sizes is not about speed at all, but about which solvers still run.

No new algorithm is claimed. The ingredient that carries the logistic result,
recomputing the Lipschitz constant on the current working set, is published:
Wang, Yu, Liang, Wu and Yu, An Iterative Reduction FISTA Algorithm for
Large-Scale LASSO, SIAM Journal on Scientific Computing 44(4), 2022, use it on
the same epsilon data at the same penalty weight and attribute their speedup
to the same smaller constant. What is ours is the measurement, the attribution
between ingredients, and the derivation of the recipe from measured symptoms.

\subsection{Conventions}
\label{app:log:conventions}

\begin{itemize}
\item Hardware. An NVIDIA RTX 4090 (sm\_89, local, shared with an LLM
 inference server) and an NVIDIA RTX PRO 6000 Blackwell Max-Q (sm\_120,
 97.9 GB, a shared four-GPU server of which one GPU is used for timing).
 CPU baselines run on an Intel Core i7-12700K, single threaded, one process
 pinned to one performance core.
\item Certificates. Always evaluated in float64 and in the original problem
 coordinates, whatever precision or rescaling the solver uses internally.
 LP and OT: the PDLP-style relative gap (primal residual, dual residual,
 objective gap). QP: OSQP's primal and dual residuals plus the duality gap.
\item Timing. Our solver's execution time excludes JIT compilation (reported
 separately). Baseline times include their own setup, which favours our
 solver on very small instances; this is stated wherever it matters.
\item Data. Every problem family is benchmarked on downloaded real data with
 a source URL and a sha256 manifest. Synthetic generators appear only in
 unit tests and early probes and are marked as synthetic below.
\item Instance manifests. Every iteration count names its data seed and draw
 size, because both change the sampled instances (Section~\ref{app:log:a1}).
\end{itemize}

\subsection{Data provenance}
\label{app:log:data}

\begin{center}
\small
\begin{tabular}{p{2.3cm}p{4.6cm}p{5.6cm}}
\toprule
Family & Data & Source \\
\midrule
TV denoising & Kermany 2017 retinal OCT B-scans (real clinical speckle) &
 Hugging Face \\
LP & 297 instances: MIPLIB 2017 (236), fctp transportation (32), Mittelmann
 network (7); 47 above $10^6$ nonzeros & MIPLIB and Mittelmann benchmark
 pages; manifest \\
Min-cost flow & road networks (DC, DE), a vision instance (246K nodes, 1.43M
 arcs), DIMACS netgen & \url{http://lime.cs.elte.hu/~kpeter/data/mcf}
 (no reader yet) \\
Elastic net & rcv1, real-sim, news20, E2006 & LIBSVM data sets \\
Convex QP & Maros-Meszaros, 138 instances, $n$ from 2 to 93263 &
 \url{https://github.com/qpsolvers/maros_meszaros_qpbenchmark} \\
Optimal transport & DOTmark 1.0: 10 classes, resolutions 32 to 512, 500
 images & \url{https://www.stochastik.math.uni-goettingen.de/DOTmark_1.0.zip} \\
\bottomrule
\end{tabular}
\end{center}

Synthetic generators used in probes: grid and random-regular min-cost flow,
point-cloud optimal transport, and multi-commodity flow built on the grid
generator.

\subsection{The factorization $T = N \times t$ on TV denoising}
\label{app:log:factorisation}

Time to a certified solution factors into the iteration count $N$ and the
per-iteration cost $t$. Measured on TV denoising of $256\times256$ OCT
patches, six images, RTX 4090, tolerance $10^{-4}$:

\begin{itemize}
\item $N$ is invariant to hardware: 584, 597 and 683 at cadence $K=16, 64,
 256$, identical across two GPU architectures, both precisions, both
 execution structures, and every kernel.
\item Scheme identity is inert at coarse cadence: PDHG and Condat-Vu give 683
 and 683 at $K=256$. Step sizes are worth $4.31\times$ on $N$ and decide
 feasibility at $10^{-6}$ (the balanced setting certifies 0 of 6).
\item The precision gene was confounded with execution structure: the
 float64 path ran a device loop with a traced inner trip count, the float32
 path ran host-chunked phases with a static one. Separated, float32 buys
 nothing inside the device loop and $1.58\times$ once the loop is chunked.
\end{itemize}

\begin{center}
\small
\begin{tabular}{llccc}
\toprule
Execution & Precision & $K=16$ & $K=64$ & $K=256$ \\
\midrule
device loop & f64 & 22.54 & 20.39 & 19.78 \\
device loop & f32 & 22.99 & 20.41 & 19.78 \\
host phases & f64 & 27.18 & 16.19 & 13.29 \\
host phases & f32 & 28.00 & 12.23 & 8.43 \\
\bottomrule
\end{tabular}

Per-iteration cost $t$ in microseconds.
\end{center}

The loop penalty is one line of code: a traced inner trip count costs
$2.02\times$ per iteration at $n=256$ and $1.18\times$ at $n=496$ against a
static one (nesting and \texttt{while\_loop} are free). On the CPU the same
choice costs $1.24\times$ at $n=128$ and inverts to $0.84\times$ at $n=256$.

\paragraph{Resolved discrepancy.}
\label{app:log:a1}
Two artifacts printed mean iteration counts 682.667 and 640.0 for the same
setting. The cause was the data draw, not a broken invariance: the harnesses
used data seeds 20260907 and 0 and draw sizes 12 and 6, and the two sampled
sets shared 0 of 6 scans. Rerun on the catalog's draw, the fusion harness
gives 683. Artifact: .

\subsection{Temporal fusion kernel}
\label{app:log:fusion}

A CUDA kernel runs $T$ PDHG or Condat-Vu iterations per launch on a
$(B+2T)\times(B+2T)$ tile with halo in shared memory. Numerical admission
against $T$ sequential reference launches: bitwise identical at $\rho=1$, at
most $4.649\times10^{-15}$ at $\rho=1.9$, on both architectures. Total wall
time over six images to a certified $10^{-4}$ gap, $K=256$, RTX 4090:

\begin{center}
\small
\begin{tabular}{rccccc}
\toprule
$n$ & baseline (s) & static loop & best fusion (tile) & total & worst tile \\
\midrule
128 & 0.06647 & $1.673\times$ & $1.571\times$ ($B=8$, $T=4$) & $2.629\times$ & $1.320\times$ \\
256 & 0.07733 & $1.347\times$ & $1.386\times$ ($B=24$, $T=4$) & $1.867\times$ & $1.298\times$ \\
384 & 0.13486 & $1.291\times$ & $1.266\times$ ($B=24$, $T=4$) & $1.634\times$ & $0.926\times$ \\
496 & 0.22564 & $1.093\times$ & $1.240\times$ ($B=32$, $T=4$) & $1.356\times$ & $0.743\times$ \\
\bottomrule
\end{tabular}
\end{center}

On the Blackwell GPU the end-to-end gains are $1.73\times$, $1.45\times$,
$1.63\times$ and $1.15\times$ at the same four sizes. The best tile agrees
across devices at $n=128$ and $256$ and disagrees at $384$ and $496$; at
$n=496$ on the Blackwell even the best tile is $0.925\times$ of the static
loop, so the right decision there is not to fuse. The wrong tile is slower
than not fusing at $n=384$ and $496$. On the catalog's own instance draw at
$n=256$, $K=256$, the gain is $1.857\times$. This remains the only result in
the project that beats software people actually use (XLA's compiled loop).

\subsection{Predicting the recipe instead of searching: two negatives}
\label{app:log:diagnostics}

Artifact: .

\begin{itemize}
\item Arithmetic intensity does not pick the fusion tile. Roofline models
 scored against the measured sweep (7 sizes, 6 tiles) reach mean Spearman
 correlation 0.355 (roofline), 0.363 (plus occupancy), 0.282 (plus wave
 quantization) and 0.567 (plus a fitted launch cost), and pick the true best
 tile in at most 1 of 7 sizes. Roofline is anti-correlated in the
 overhead-bound regime ($-0.943$ at $n=128$) and correct only once the
 kernel is bandwidth bound ($+0.829$ at $n\ge992$).
\item Early sharpness of the gap curve does not predict $N$. The pooled
 correlation across families reversed within families (a Simpson's paradox)
 and a leakage defect was found in the first version of the probe; sharpness
 is not used as a predictor.
\end{itemize}

\subsection{Problem families beyond TV}
\label{app:log:families}

\paragraph{Min-cost network flow (synthetic grid and regular generators).}

The execution-structure gain replicates on a sparse LP: mean $1.386\times$
(range 1.277 to 1.492, 24 runs), with $N$ identical across loop structures
in 24 of 24 runs and identical between PDHG and Condat-Vu. The step sizes
that were merely slow on TV are inadmissible here: the construction gate
rejects them before execution ($t s \|L\|^2 = 1.02267$), 24 of 48
configurations. In the structure sweep (2106 runs, 1755 admitted, 351
rejected) admission depends only on the fill $\theta = t s \|L\|^2$: every
$\theta\le0.99$ is admitted and $\theta=1.02$ never is. The admissible
region is family dependent: $\|A\|=\sqrt{m+n}$ for transport, a region
shrinking as $\Theta(1/K)$ in the number of commodities for multi-commodity
flow, and a norm that saturates at 4 on grids.

\paragraph{The primal-weight rule.}
 A step-ratio rule
$r = C/\omega_0^p$ with $\omega_0 = \|c\|/\|b\|$ was fitted from data. It ties
PDLP's primal-weight heuristic ($p=2$): $1.0704\times$ against $1.0704\times$
of the per-instance oracle on 88 ultrawide netflow instances, $1.1822\times$
against $1.1829\times$ on 80 fine-grid instances, $1.1984\times$ against
$1.2020\times$ on 108 wide instances. The fitted exponent (about 1.8) moved
from 2.00 to 1.80 as grid censoring was removed. Transferred to transport and
multi-commodity flow, the rule reaches $1.125\times$ of the oracle on 162
held-out instances against $1.825\times$ for the best fixed ratio. We
reproduce PDLP's rule; we do not beat it.

\paragraph{Elastic net (LIBSVM).}
Artifact: .
Coordinate descent with working sets (celer, adelie, glmnet) dominates. A
full-gradient method loses $13.6\times$ (rcv1) and $18.8\times$ (news20) core
for core. On the RTX 4090 against one CPU core it ties rcv1 (0.130 s against
0.137 s), wins news20 by $1.45\times$, and loses about $550\times$ on E2006,
which must be centered (uncentered, one near-constant feature solves it at
every regularization value tried). Verdict: a poor target, because the
in-class state of the art is not first-order.

\subsection{Convex QP against author code (Maros-Meszaros)}
\label{app:log:qp}

 This is the first comparison
against solvers written by their authors. Every solver's returned primal and
dual pair is scored by the same certificate at $10^{-4}$; a baseline whose
own stopping rule returns a point failing it is rerun with its tolerance
tightened tenfold, up to three times. OSQP and SCS are first-order (ADMM) and
in our class; Clarabel is an interior-point method, reported for context.
Our solver is Condat-Vu (PDHG is rejected by the gate, since the quadratic
has no cheap prox), float64, cadence 64, capped at $5\times10^4$ iterations,
with and without the Ruiz rescaling gene, which was extended to QP for this
comparison (KKT equilibration as in OSQP).

\begin{center}
\small
\begin{tabular}{lcc}
\toprule
Solver & solved (of 137) & time where both solve, vs OSQP \\
\midrule
Clarabel & 132 & \\
OSQP & 123 & \\
SCS & 123 & \\
ours, plain & 55 & $4.47\times$ slower (geometric mean, 55 instances) \\
ours, Ruiz(10) & 63 & $5.80\times$ slower (geometric mean, 63 instances) \\
\bottomrule
\end{tabular}
\end{center}

The 14 instances where our solver appears faster than OSQP are, with one
exception, tiny (the HS family and similar), where our timer excludes
compilation and OSQP's includes its factorization: timing artefacts. The
one large instance, BOYD1, where plain Condat-Vu passed the certificate
faster than OSQP (33.2 s against 182.9 s), is not a solution: its point
violates rows by up to 9.7 times their own bound and its objective is 84\%
off (accuracy audit below). It was reported as a win before the audit; that
claim is retracted. On CONT-200, CONT-201 and
CONT-300 OSQP and SCS need 900 to 3700 s at the tightened tolerance and
Clarabel 6 to 8 s, while our solver does not certify within its iteration
cap. Ruiz rescaling rescues some
instances (DUALC1) and breaks others (CVXQP1\_S), because it moves the
implied primal-dual step ratio. What OSQP has and we do not: an adaptive step
ratio, factorization-based steps, and solution polishing.

\paragraph{Accuracy audit: what a certificate pass means.}
 The KKT certificate
follows OSQP and divides every residual by one global scale, so on badly
scaled instances it admits points that are not solutions. Every pass was
therefore re-judged, for our arms and for the baselines rerun at the exact
tolerance each needed. Strict: every row violated by at most $10^{-4}$ of its
own bound and the objective within $10^{-3}$ of Clarabel run to $10^{-9}$
(trusted where it reports solved, 132 instances). Borderline: violations up
to $10^{-2}$ with the objective within $10^{-2}$, approximate answers whose
status depends on the normalization. Gross: a violation above $10^{-2}$ or an
objective more than $10^{-2}$ off.

\begin{center}
\small
\begin{tabular}{lrrrrr}
\toprule
Solver & passes judged & strict & borderline & gross & not grossly wrong \\
\midrule
Clarabel & 132 & 109 & 14 & 9 & 123 \\
OSQP & 119 & 51 & 39 & 29 & 90 \\
ours, Condat-Vu + Anderson & 66 & 36 & 19 & 11 & 55 \\
ours, Condat-Vu & 55 & 24 & 21 & 10 & 45 \\
\bottomrule
\end{tabular}
\end{center}

Every solver was stopped when this style of certificate reached $10^{-4}$, so
the table describes what the certificate admits, not what each solver can
reach with a stricter stop. OSQP, whose own residuals the certificate copies,
has 29 gross failures (QGROW15, for example, violates a row by 87 times its
bound). Clarabel's 9 are objective errors on S268, HS268, GOULDQP3 and the
LISWET family, where a violation of $2\times10^{-4}$ moves the objective by 4 to
90\% for every solver. The ranking survives the audit, and so does the
Anderson gain (10 more answers not grossly wrong, 12 more strict; 8 and 10
after the composer fixes below); the
certificate-only counts reported above (OSQP 123, ours 55 and 63) overstate
every solver's accuracy. OSQP's four slow passes (BOYD1, CONT-200, CONT-201,
CONT-300) were not re-judged. All later QP claims use this judge.

\paragraph{Correction and rerun.}
The last column of the table above first printed 116, 89, 54 and 44, which is
not strict plus borderline in any row; recounting from
 with the tier definitions above gives 123,
90, 55 and 45, now in the table. After the two composer defects of
the timing correction were fixed, both of our arms were rerun on all 137
instances and re-judged . Condat-Vu is unchanged (55 passes;
24 strict, 21 borderline, 10 gross). Condat-Vu + Anderson passes 64 instead of
66, losing DTOC3 and PRIMALC2 (dual infeasible as distributed), with 34 strict,
19 borderline and 11 gross: 53 against 45 answers are not grossly wrong, an
Anderson gain of 8, and 10 more strict.

\subsection{A knowledge base of classical schemes}
\label{app:log:kb}

Artifact: . Methods adopted in LP practice
since 2021 (restarts, Halpern reflection, averaging, primal weights) trace
back to fixed-point papers from 1950 to 1980 through the monotone-inclusion
literature. The knowledge base records, for each classical scheme, its
original paper, update, side condition, guarantee, genome location, gate
rule, and whether it has already been used in general primal-dual splitting
and in LP or QP practice (with evidence or the queries that found nothing).

\begin{center}
\small
\begin{tabular}{p{2.6cm}p{2.2cm}p{7.6cm}}
\toprule
Scheme & Origin & Status and outcome \\
\midrule
Krasnoselskii-Mann, Halpern, restarts & 1953 to 1967 & implemented (genome wrapper genes) \\
Anderson acceleration & 1965 & already used for PDHG on LP (arXiv 2508.08062); not frontier \\
Passty averaging & 1979 & already used (restarted averaged PDHG, arXiv 2206.03041); not frontier \\
Haugazeau & 1968 & implemented as an outer wrapper, with the projection evaluated in the PDHG metric; admission-tested. Rejected by measurement on QP: solves 8 and 9 instances against 55 and 63 without it, needs $20.5\times$ and $34.3\times$ more iterations where both solve, and solves nothing new \\
Bregman (entropic) & 1967 & implemented as an entropic PDHG. The gate admits it only where the constraints imply a simplex (transport), with a size-independent step condition $t s \|A\|_{1\to2}^2<1$, $\|A\|_{1\to2}=\sqrt2$. Needs a dual recovery step; never certified on DOTmark \\
Solodov-Svaiter & 1999 & frontier, not yet tried \\
Ishikawa & 1974 & not found in primal-dual practice; not tried \\
\bottomrule
\end{tabular}
\end{center}

Features harvested from author code and recorded for implementation as
combinable genes: an adaptive step ratio (OSQP's adaptive $\rho$, PDLP's
primal weight), a factorization-based prox for the quadratic (which would
make PDHG admissible on QP), and polishing (an active-set reduced solve
accepted only if the certificate improves).

\subsection{Optimal transport on DOTmark}
\label{app:log:ot}

 Baseline: POT's exact
network simplex (\texttt{ot.emd}) on one CPU core, 0.04 to 0.85 s per pair at
$32\times32$ (about $10^6$ arcs) and 0.8 to 44 s at $64\times64$ (about
$1.7\times10^7$ arcs).

\begin{itemize}
\item Execution structure. Writing the transport operator as row and column
 sums of the plan instead of scatter-adds over arcs (same arithmetic)
 cuts the RTX 4090's per-iteration time from 243 to 75 microseconds at
 $32\times32$ and from 6.32 to 1.07 ms at $64\times64$. One CPU core needs
 15 ms per iteration at $32\times32$ even with the faster form, so the CPU
 arm was not pursued.
\item Step ratio. A fixed ratio grid over $[0.03, 30]$, and a wider one, certified
 nothing in $5\times10^4$ iterations: the plan stayed at zero. With the ratio
 set from $\omega_0 = \|c\|/\|b\| \approx 8.8\times10^6$, plain PDHG certifies
 $10^{-4}$ in 5184 iterations. Reflected Halpern restarts did not help; the
 entropic scheme never certified.
\item Certificate weakness. The LP certificate normalizes the primal residual
 by $1+\|b\|$; for probability measures $\|b\|\approx0.01$, so a $10^{-4}$
 pass is almost an absolute bound and admitted a plan with $4.7\%$
 objective error. Transport results are therefore judged by relative error
 against the exact cost.
\item Head to head. At $32\times32$ and $10^{-6}$ (objective error
 $5\times10^{-7}$): 3.0 s on the GPU against 0.21 s for POT, $14\times$
 slower. At $64\times64$ (ClassicImages): 6.2 s against 19.4 s, but with
 $4.7\%$ objective error, which is not a solution; 20.4 s against 19.4 s at
 $7\times10^{-4}$ objective error, a tie. No win at matched meaningful
 accuracy.
\end{itemize}

\paragraph{Creation loop on transport (corrected composer).}
On CauchyDensity at resolution 32, pair 1001-1002 (tolerance $10^{-6}$, primal
weight $\omega_0/30$, one Blackwell GPU), the gates admit 19 of 48 compositions
of PDHG. Plain PDHG certifies in 69,248 evaluations and 4.27 s. The best
composition is reflected Halpern with a fixed restart every 512 steps: 61,440
evaluations and 3.90 s, $1.13\times$ fewer evaluations and $1.10\times$ faster.
With a restart every 64 steps it is roughly neutral; residual-triggered restarts
cost more evaluations or fail, and unanchored Halpern and Anderson acceleration
reach the cap of $10^5$ evaluations. A first run of this loop reported the
residual-triggered recipes as $12\times$ faster than plain PDHG; that was an
execution artifact of the composer, not a
property of the recipes.

\subsection{Certified adaptive step ratio}
\label{app:log:adaptive}

PDHG and Condat-Vu are proximal steps in the metric
$M = [[I/t, -L^\top], [-L, I/s]]$, so changing the primal weight
$\omega = \sqrt{s/t}$ changes the metric. With $\theta = \sqrt{ts}\,\|L\| < 1$,
$M \succeq (1-\theta)D$ for its block diagonal $D$, so a move from $(t,s)$ to
$(t',s')$ satisfies $M' \preceq (1+\eta)M$ with
$\eta = \max(0, t/t'-1, s/s'-1)/(1-\theta)$ (checked numerically over 200
random moves). The controller applies PDLP's primal-weight update every 512
iterations, places every $(t,s)$ on the gate boundary, and stops adapting
once $\sum\eta$ reaches a declared budget, which keeps the variable metric
convergence theorem of Combettes and V\~u in force. The comparison arm runs
through the same runner with the controller off.

\paragraph{QP (Maros-Meszaros, 137 instances, CPU, tolerance $10^{-4}$).}
With Ruiz rescaling the adaptive ratio passes the certificate on 73
instances against 65 with the fixed ratio (15 gained, 7 lost); where both
pass it needs $0.53\times$ the iterations (geometric mean; fewer on 35 of 58).
Without rescaling: 62 against 55 (9 gained, 2 lost), $0.50\times$ the
iterations. The budget froze 8 runs; the median run made 30 updates. Under the
accuracy audit (Section~\ref{app:log:qp}) the gain survives: with Ruiz
rescaling 59 against 50 answers are not grossly wrong (33 against 25 strict),
and without rescaling 53 (25 strict). That makes it our best QP arm, ahead of
Condat-Vu with Anderson (54), but OSQP has 89 such answers and Clarabel 116,
and where both pass the certificate it remains $7.8\times$ slower than OSQP. A
real gain inside the method, not a win.

\paragraph{Optimal transport (DOTmark $32\times32$, 31 pairs, RTX 4090, tolerance $10^{-6}$).}
Started at PDLP's initial weight $\omega_0 = \|c\|/\|b\|$, neither the fixed nor
the adaptive ratio certifies any pair within $10^5$ iterations, although
adaptation cuts the median objective error from $2.7\times10^{-2}$ to
$1.1\times10^{-3}$; it moves the ratio at every epoch and never settles. The
fixed ratio $\omega_0/30$ calibrated on a held-out pair certifies 29 of 31, a
median $19\times$ slower than POT (best $4.7\times$). At $64\times64$ nothing
certifies within $10^5$ iterations (105 s against 39 to 47 s for POT). On
transport, adaptation does not replace calibration.
Adapting for a bounded number of updates and then freezing (epochs of 2048
iterations) does not rescue it: at most 10 updates from $\omega_0$ certifies 5
of 31 pairs, at most 30 updates certifies 12, and 10 updates started from the
calibrated ratio certifies 6, against 29 for the calibrated fixed ratio. The
runs that do certify are faster (a median 11.7 to $14\times$ POT against
$21.6\times$), so the rule helps when it lands, but it does not land reliably.

\subsection{Elastic-net logistic regression}
\label{app:log:logreg}

Proposed by N.~Devanathan: elastic-net logistic regression on epsilon (dense,
$4\times10^5\times2000$), rcv1 (sparse), Tabula Sapiens (1.1M cells) and
ogbg-molpcba (437,929 molecules, 128 assays, multitask), against cuML's GPU
quasi-Newton (OWL-QN), skglm, and glmnet-style adelie. Two items to add in
order: fixed diagonal preconditioning followed by adaptation, and
certificate-driven structure reduction by safe screening. The possibility
the user singled out is that screening shrinks the problem mid-solve until
the best device flips (GPU first, CPU afterwards), which no existing solver
does automatically.

Built so far. Atoms for the logistic loss ($L_f = \|X\|_2^2/(4n)$) and the
elastic-net penalty (scaled soft threshold). The certificate is the relative
duality gap at the dual point $u = \nabla F(Xw)$; for $\alpha<1$ the penalty's
conjugate is smooth, for $\alpha=1$ the point is rescaled into dual
feasibility. FISTA with O'Donoghue--Cand\`es gradient restart is a new scheme
with gate $0 < a \le 1/L$ and an independent reference in the equivalence
battery. Gap Safe screening (Ndiaye, Fercoq, Gramfort and Salmon, 2017)
uses the radius $\sqrt{\mathrm{gap}/(2n)}$, since the dual is $4n$-strongly
concave (the binary-entropy term has curvature at least 4).

Defect caught. For $\alpha=1$ the bare strict screening test discarded two
active features at an exactly converged point: active features sit on
$|x_j^\top u^\star| = \lambda$, and the rescaling into dual feasibility can leave
them one rounding step below $\lambda$. A relative margin of $10^{-9}$ fixes it;
the safety test (no screened feature nonzero at a $10^{-12}$-accurate optimum)
passes for $\alpha \in \{0.5, 0.9, 1\}$.

Baseline check. On synthetic data every baseline certifies the same
objective under our gap: cuML with $C = 1/(n\lambda)$ reaches $2\times10^{-15}$,
confirming the mapping; skglm, scikit-learn saga and adelie also pass.
Baselines run in isolated environments on both machines. The training file
of epsilon is being resumed after a truncated first download.

\paragraph{First head-to-head (epsilon test split, real data).}
$n=10^5$, $p=2000$, $\alpha=0.5$, $\lambda=0.1\lambda_{\max}$, target relative duality
gap $10^{-6}$ on our certificate; GPU runs on the RTX 4090, CPU baselines on
one core. 

\begin{center}
\small
\begin{tabular}{llrrl}
\toprule
Solver & device & seconds & iterations & gap \\
\midrule
cuML (OWL-QN, float32) & 4090 & 0.183 & 35 & $2.2\times10^{-7}$ \\
ours, FISTA (float64) & 4090 & 0.668 & 192 & $6.3\times10^{-9}$ \\
ours, FISTA (float32 iterates, float64 certificate) & 4090 & 0.722 & 192 & $5.9\times10^{-9}$ \\
ours, forward-backward & 4090 & 4.88 & 1408 & $7.0\times10^{-7}$ \\
adelie (30-point path) & 1 core & 9.05 & & $2.2\times10^{-8}$ \\
scikit-learn saga & 1 core & 21.8 & & $8.5\times10^{-11}$ \\
skglm (tolerance tightened to $10^{-6}$) & 1 core & 59.2 & & $3.5\times10^{-8}$ \\
\bottomrule
\end{tabular}
\end{center}

FISTA on the GPU is $13.5\times$ faster than the fastest single-core baseline,
which is not a like-for-like hardware comparison. On the same GPU, cuML is
$3.7\times$ faster than our FISTA. Float32 iterates buy nothing (0.722 s against
0.668 s). The iteration is expected to be memory bound, since the design is
1.6 GB in float64 and is read twice per step (for $Xw$ and for $X^\top(\cdot)$);
a fused single-pass gradient kernel that reads each row block once for both
products is the natural next step. The first cuML attempt ran out of GPU
memory because our process still held the data next to the language-model
server; it was rerun alone.

\paragraph{Diagnosis: where the step's time goes (RTX 4090, epsilon test split).}
 In float64 the
step is memory bound at about 90\% of peak bandwidth: $Xw$ takes 1.74 ms (918
GB/s), $X^\top v$ 1.84 ms (871 GB/s), and the gradient 3.52 ms, two full reads
of $X$, which is essentially the whole measured FISTA iteration (3.48 ms). In
float32, $Xw$ is bandwidth bound and twice as fast (0.89 ms), but $X^\top v$
runs at only 298 GB/s (2.70 ms), so the float32 gradient is no faster than
the float64 one. Storing $X^\top$ explicitly, or writing the product as $vX$,
changes nothing (294 to 296 GB/s): the loss is the library kernel XLA selects
for this shape (2000 outputs, each a sum of $10^5$ terms), not the memory
layout. Mixed precision therefore pays only with a kernel that reaches full
bandwidth on the transposed product; a fused single pass in float32 would
cost about 0.9 ms per gradient, four times less than float64, and bfloat16
storage would halve that again. Low-precision data changes the problem, so
reaching a $10^{-6}$ gap on the true problem needs a precision schedule that
returns to float64 gradients as the gap closes; the certificate stays float64
throughout.

\paragraph{The predicted remedy was wrong, and no kernel was needed.}
The paragraph above proposes a fused single-pass gradient kernel. No such kernel
can exist: the transposed product consumes a vector that depends on the result
of the forward product, so no single pass over the design can produce both. The
diagnosis of the cause was also wrong. Layout is not the problem, as that
paragraph already half-observed: $vX$, an einsum, and an explicitly materialized
transpose all land within a per cent of each other. But the same float32 data
accumulated in float64 runs twice as fast, which places the loss in the
accumulation over the length $10^5$ reduction rather than in the traffic or the
pass structure. Four ways of writing the same product, measured on the $10^5
\times 2000$ design:

\begin{center}
\small
\begin{tabular}{lrrl}
\toprule
Form & ms & GB/s & relative error \\
\midrule
$X^\top v$, as shipped & 2.589 & 309 & $2.9\times10^{-4}$ \\
contracting with a float64 accumulator & 1.417 & 564 & $1.0\times10^{-14}$ \\
reduction split into 16 chunks & 0.844 & 948 & $2.9\times10^{-4}$ \\
row sum of the elementwise product & 0.858 & 932 & $1.2\times10^{-7}$ \\
\bottomrule
\end{tabular}
\end{center}

The form we shipped is dominated on both axes at once: it is three times slower
than the third row and no more accurate. The last row is what the solver now
uses in float32. It reaches the same bandwidth as the forward product, it
materializes no intermediate (0.858 ms for 0.80 GB; materializing would cost
about 2.4 ms), it places no divisibility condition on the sample count, and it
is more accurate than the alternative of equal speed by a factor of $2400$,
which matters because the certificate is float64. Float64 keeps the direct form,
being already at $90\%$ of peak.

End to end on the working-set configuration at $\alpha=0.5$ and
$\lambda=0.1\lambda_{\max}$, float32 falls from $0.0564$ s to $0.0467$ s, a
factor of $1.21$, at an unchanged 64 evaluations and an equivalent certified
gap, so the gain is entirely in cost per iteration. Float32 now beats float64
($0.0606$ s) by $1.30\times$, where before it did not, which revises the earlier
row of this section reporting that float32 iterates buy nothing: that row is a
full-width plain FISTA configuration and it stands as measured, but the
conclusion drawn from it does not generalize. The realized $1.21\times$ is
smaller than the $3.07\times$ on the isolated product because the solver runs on
a working set of 125 to 500 columns rather than all 2000. The revised prediction
for a saturated float32 gradient is a factor of two over float64, not the factor
of four predicted above, because two passes over the design are unavoidable.
This is one device. The second machine could not be used to check it, its CUDA
installation being broken pending a reboot.

\paragraph{Full training set (epsilon, $4\times10^5$ samples, Blackwell GPU).}
The training split does not fit next to the language-model server on the
RTX 4090 in float64, so it ran on one RTX PRO 6000 Blackwell GPU with the same
settings ($\alpha=0.5$, $\lambda=0.1\lambda_{\max}$, target gap $10^{-6}$ on our
certificate); CPU baselines on one core. 

\begin{center}
\small
\begin{tabular}{llrrl}
\toprule
Solver & device & seconds & iterations & gap \\
\midrule
cuML (OWL-QN, float32) & Blackwell & 0.755 & 37 & $3.0\times10^{-7}$ \\
ours, FISTA (float64) & Blackwell & 1.53 & 192 & $6.7\times10^{-9}$ \\
ours, forward-backward & Blackwell & 11.2 & 1408 & $7.3\times10^{-7}$ \\
adelie (30-point path) & 1 core & 39.0 & & $5.4\times10^{-8}$ \\
scikit-learn saga & 1 core & 99.2 & & $1.7\times10^{-10}$ \\
skglm (tolerance tightened to $10^{-6}$) & 1 core & 247.9 & & $8.5\times10^{-9}$ \\
\bottomrule
\end{tabular}
\end{center}

On the same GPU cuML is $2.0\times$ faster than our FISTA (against $3.7\times$ on
the test split on the 4090). Our FISTA is $25\times$ faster than the fastest
single-core baseline, which is not a like-for-like hardware comparison. The
logistic certificate is the relative duality gap, which a diverging iterate
cannot pass, so these passes do not share the weakness found in the QP
certificate.

\paragraph{Correction: the step size was not timed.}
FISTA and forward-backward need a Lipschitz constant $L$ of the gradient. In
every row above, and in the creation-loop rows of
Section~\ref{app:log:kbloop}, it was computed before the timer started, by 100
power iterations on the host (200 passes over $X$), while each baseline's time
includes all of its own setup. With $L$ computed on the GPU inside the timed
region (30 power iterations), FISTA takes 0.775 s instead of 0.671 s on the
4090 test split and 1.775 s instead of 1.536 s on the Blackwell training split.
Every row from here on times the step size.

\subsection{The knowledge-base creation loop}
\label{app:log:kbloop}

 Restarted reflected Halpern PDHG was reached
by transferring general operator ideas (Halpern's anchor, then restarts, then
reflection) to one averaged operator, PDHG. The loop automates that route for
another averaged operator, forward-backward ($\alpha = 2/(4-aL)$ for $a<2/L$),
on elastic-net logistic regression. Ingredients come from the knowledge base:
Krasnoselskii-Mann relaxation, Halpern anchoring, fixed and residual-triggered
restarts, reflection under the anchor, safeguarded Anderson acceleration, and
FISTA momentum as a separate base. Two gate rules were added for this:
residual-triggered restart of the anchor (sound because Halpern's residual
bound forces every segment to end) and Anderson acceleration with the
Zhang, O'Donoghue and Boyd safeguard.

The composer enumerates 97 compositions; the gates reject 58 before anything
runs (for example relaxation beyond $1/\alpha$ without an anchor, momentum
combined with an anchor, two outer extrapolations together) and admit 39,
each run to the certified duality gap. Plain and anchored recipes reproduce
the existing runners exactly (same evaluation counts and iterates), which
validates the composer.

Novelty is checked against the knowledge base's evidence for sparse
generalized-linear-model practice, and the check mostly came back negative:
Anderson acceleration of proximal gradient on logistic regression is
published (Mai and Johansson, ICML 2020), a Halpern anchor has been used on
LASSO (Numerical Algorithms, 2026), and restarting FISTA on such problems is
standard. The only ingredient recorded as unused in that practice, on two
searches, is a periodic restart of a Halpern anchor on forward-backward. A
related June 2026 preprint applies restarted reflected Halpern acceleration
to augmented primal-dual methods (arXiv 2606.16552).

The composer is generic over base operators: every wrapper acts on the
operator's state as a pytree of arrays, so it applies unchanged to forward-backward, FISTA,
PDHG, Condat-Vu, Douglas-Rachford and Davis-Yin, and it reproduces the policy
runner exactly on forward-backward, on Condat-Vu with a restarted anchor for a
real QP, and on PDHG for a transport LP. On the QP instance AUG3DCQP the gates
choose Condat-Vu as the base (19 admitted, 29 rejected) and the loop transfers
Anderson acceleration to it without any QP-specific code: 257 evaluations and
0.030 s against 6400 and 0.179 s for plain Condat-Vu, to a certified $10^{-6}$
KKT residual (both times were measured before the flattening defect of
the timing correction was fixed; the evaluation counts are unaffected). The restarted Halpern chain behind restarted Halpern PDLP loses
there (12,500 to 19,000 evaluations), so that chain is not universally
beneficial. Anderson is recorded as already used in QP practice (SCS uses it),
so this is an automatic rediscovery, not a new algorithm. Family-wide, Condat-Vu + Anderson passes the
certificate on 66 of 137 instances against 55 for plain Condat-Vu; under the
accuracy audit (Section~\ref{app:log:qp}) 55 against 45 are not grossly
wrong; after the composer fixes the rerun gives 64 against 55 passes and 53
against 45 not grossly wrong (Section~\ref{app:log:qp}).

On the epsilon test split at a $10^{-8}$ gap (RTX 4090, after the Anderson accounting fix)
the gates admit 39 of 97 compositions; 33 reach the certificate and 6 hit the cap of
20,000 evaluations, all six unrestarted Halpern anchors, whose $O(1/N)$ rate is too slow
for this accuracy. Plain FISTA is fastest (192 evaluations, 0.947 s). The best
composition on forward-backward uses the step $1.9/L$ (1088 evaluations, 3.90 s, with or
without relaxation); Anderson acceleration lowers the evaluations to 641 to 705 but not
the time (4.8 to 5.0 s); restarted Halpern anchors need 2176 to 2304 evaluations (7.7 to
8.1 s). No composition of these ingredients beats FISTA on this instance: the transfer
route that produced restarted Halpern PDHG finds nothing here.

For the proof-check stage, the performance-estimation package PEPit is
installed in an isolated environment and reproduces the Halpern worst case
exactly ($\|x_N - T x_N\|^2 \le 0.0330579\,\|x_0 - x^\star\|^2$ at $N = 10$); its
fixed-point, monotone-inclusion and composite families are the starting point
for checking a transferred composition before it runs.

That check is now part of the loop .
A fixed-coefficient recipe on an $\alpha$-averaged base, $T = (1-\alpha)I + \alpha N$
with $N$ nonexpansive, relaxed by $\rho$ and iterated plainly or with the
Halpern anchor, becomes a performance-estimation problem whose value is the
worst case of $\|x_N - T x_N\|^2$ over the whole operator class. It reproduces
Halpern's bound $(2/11)^2$ at $N=10$, and it agrees with the gates without being
told what they decided: reflection without an anchor ($\rho\alpha = 1$), which the
gate rejects, shows no worst-case decay from $N=10$ to $N=20$, while strict
relaxation and reflection under the anchor, which the gate admits, contract.
The gate and this analysis are two independent derivations of the same
admissibility, one symbolic and one numerical. Restarts and Anderson
acceleration have iterate-dependent coefficients, so they are not
fixed-coefficient problems and are left to the certified runs.

\subsection{Diagnose, mutate, re-run: closing the gap to cuML}
\label{app:log:diagnose}

The composition loop of Section~\ref{app:log:kbloop} found nothing that beats
FISTA on epsilon, so a diagnosis stage was added in front of it. It splits the
gap to the baseline into the two factors of $T = N \times t$ and turns
measurements into symptoms, each mapped to knowledge-base ingredients that act
on one factor:
\begin{itemize}
\item the solution is sparse (197 of 2000 features) but every step reads all
 columns: Gap Safe screening with mid-solve re-planning (acts on $t$);
\item the step is memory bound (about 90\% of peak bandwidth in float64):
 float32 storage, or a fused gradient kernel where the float32 transposed
 product is slow, as on the 4090 but not on the Blackwell (acts on $t$);
\item the baseline needs far fewer iterations (37 against 192): variable
 metric, Anderson acceleration, a quasi-Newton metric (acts on $N$). Anderson
 and the Halpern anchor are now reachable, by routing the configuration to the
 composition engine of Section~\ref{app:log:kbloop}, and both were tried and
 beaten on measurement; the quasi-Newton metric and the fused kernel remain
 unimplemented.
\end{itemize}
Further rules compare a run with the run it was mutated from: a re-plan that
restarts FISTA's momentum (keep it); exact-size re-plans that compile a new
shape inside the solve (pad to halving buckets compiled from $(n, p)$ alone);
iteration counts rounded up to the certificate cadence (halve $K$); safe
screening that cannot act until the gap is small (start from a working set
ranked by dual scores and let the full-problem certificate add violators, the
approach of celer and skglm); and, on a working set, the exact Lipschitz
constant of the selected columns in place of power iteration. The search is a
beam of width two over single mutations, so a step that pays only after the
next one (screening, before its re-planning costs are fixed) is not discarded;
a first greedy version stopped at 0.500 s for exactly that reason. The
reported chain is the ancestry of the best configuration. Every certificate is
the full-problem duality gap, so no mutation changes what is certified, only how
fast it is reached.

Epsilon test split on the RTX 4090, target gap $10^{-6}$, medians of five
solves after an untimed warm-up, step size timed:

\begin{center}
\small
\begin{tabular}{lrrl}
\toprule
Configuration & seconds & evaluations & gap \\
\midrule
FISTA, $K=64$ & 0.775 & 192 & $6.3\times10^{-9}$ \\
FISTA, $K=16$ & 0.642 & 144 & $9.0\times10^{-7}$ \\
\quad + Gap Safe screening, momentum kept, buckets & 0.456 & 96 & $5.5\times10^{-7}$ \\
\quad + working set from $p/16$ columns & 0.081 & 48 & $1.0\times10^{-9}$ \\
\quad + certificate every 32 steps, exact step on the working set & 0.067 & 64 & $9.4\times10^{-7}$ \\
cuML (OWL-QN, float32), same GPU, eight runs & 0.166 to 0.295 & 35 to 39 & $1.6$ to $4.6\times10^{-7}$ \\
\bottomrule
\end{tabular}
\end{center}

The median of the last five cuML runs on the 4090 is 0.227 s; one of them
needed a tolerance of $10^{-7}$ to reach the target. On the Blackwell training
split (one RTX PRO 6000 GPU; cuML alone on an identical GPU, and five times on
the same GPU) the rows read: FISTA 1.775 s, screening 1.026 s (96 evaluations),
working set 0.176 s (48 evaluations, gap $9.5\times10^{-10}$), and 0.142 s with
the certificate every 32 steps and the exact step. cuML took 0.482 to 0.636 s
after a cold first run of 1.32 s (median 0.615 s, 34 to 43 iterations), and two
of the five runs needed a tolerance tighter than $10^{-4}$ to reach the target.
So with nothing excluded but compilation, the best configuration found by hand
is $2.5\times$ faster than cuML's fastest run on the 4090 and $3.4\times$ on the
Blackwell ($3.4\times$ and $4.3\times$ against cuML's medians); the loop's own best
configurations, below, are $3.0\times$ and $3.7\times$ faster than cuML's fastest runs
($4.1\times$ and $4.7\times$ against the medians).

Where the speed comes from: the working set changes both factors. Early steps
read 125 to 250 columns instead of 2000, which is $t$; and the Lipschitz
constant of those columns is 0.0017 to 0.0031 against 0.089 for all of $X$ (a
strong common direction dominates the full spectrum), so the step is 30 to 50
times larger and the certificate is reached in 48 evaluations instead of 144,
which is $N$. Safe screening alone reaches only 96 evaluations because it
cannot act before evaluation 80.

What the loop found by itself: on the 4090, in its first version (with $L$
still computed outside the timer), the beam search from plain FISTA at $K=64$
reached $K=32$ with screening, a working set, buckets, float32 storage and
kept momentum, 0.127 s against cuML's 0.166 s. On the Blackwell, with the step
size timed, it reached the working set in its third round (0.333 s, then 0.232 s
at $K=32$, against 0.571 s). With the step size timed, the 4090 loop took screening, the working set,
$K=32$ and float32 storage, 0.102 s ($1.6\times$ cuML's fastest run), and the
Blackwell loop stopped at 0.232 s; both fell short of the best configuration
found by hand. The shortfall was traced to a gene no rule could move, the size
of the first working set, which held 250 columns for a solution with 197
nonzeros. A rule was added that halves it when it already holds the final
support. The Blackwell rerun then went, unaided, from 1.777 s through screening
(1.518 s), the working set (0.333 s), the smaller first set (0.191 s), $K=32$
(0.141 s) and float32 storage (0.132 s): ahead of the configuration found by
hand, and $3.7\times$ faster than cuML's fastest run on the same GPU model. The
4090 rerun with the new rule went from 0.777 s through screening (0.713 s), the
working set (0.148 s), the smaller first set (0.086 s), $K=32$ (0.063 s), float32
storage (0.059 s) and kept momentum (0.056 s), also ahead of the configuration
found by hand (0.067 s): $3.0\times$ faster than cuML's fastest run on the same GPU
and $4.1\times$ faster than its median. This is the harness evolving from
its own failures, the pattern CAKE (arXiv 2608.12629) applies to kernel agents,
where recurring failures become verifier rules.

\paragraph{Guided against blind mutation (Blackwell, epsilon training split).}
 Both arms search the same eight genes
(certificate cadence $K$, screening, momentum on re-planning, buckets, working
set, size of the first working set, step rule, precision) from the same start
(FISTA, $K=64$, power-iteration step), with the same beam of two, the same
certificate and the same budget of 24 evaluated configurations. The guided arm
takes its mutations from the diagnosis rules; the blind arm takes four random
single-gene changes per expanded configuration and sees only certified wall
time. Three seeds, one per GPU; medians [min, max] of the best certified time
(cuML 0.571 s): after 8 configurations, guided 0.802 s [0.801, 0.805] and blind
1.207 s [0.910, 1.239]; after 16, 0.190 s [0.186, 0.192] and 0.192 s [0.140,
0.782]; after 24, 0.140 s [0.136, 0.145] and 0.141 s [0.140, 0.248].

On this small space the diagnosis does not reach a better endpoint: one gene
(the working set) carries most of the gain, and random single-gene changes find
it within 24 configurations in two of three seeds. What the diagnosis buys here
is reliability and early progress: its three runs agree to within 7\%, while the
blind runs span 0.14 to 0.78 s at 16 configurations and one ends at 0.248 s. The
stronger evidence for the diagnosis lies before this ablation: the timing-only
composition loop of Section~\ref{app:log:kbloop} searched a space without
screening or working sets and found nothing that beats FISTA, and those
ingredients entered the search space only because the diagnosis named the
sparse solution and the gap-limited screening as the bottleneck. On the RTX 4090
(epsilon test split, cuML 0.166 s, three seeds run in sequence) the pattern
repeats: after 16 configurations guided 0.087 s [0.086, 0.087] and blind
0.150 s [0.065, 0.469]; after 24, guided 0.064 s in all three seeds and blind
0.065 s [0.064, 0.119]. On both devices the diagnosis gives the same endpoint
as a lucky blind run, every time; the blind search gives it in two runs of three.

Does the diagnosis matter more when the space a blind search must cover is
larger? The blind arm was rerun on the 4090 with a larger space of real solver
settings (six cadences and six first-set sizes instead of four, plus the number
of power iterations, the screening threshold, and FISTA against
forward-backward); the guided arm is unchanged, because its rules never propose
those genes. Three seeds, best certified time after 16 configurations: blind
0.149 s [0.135, 0.651] against 0.150 s on the small space; after 24, blind 0.087 s
[0.063, 0.119] against 0.065 s on the small space and 0.064 s for the guided arm.
The blind median gets worse as the space grows ($1.36\times$ the guided time after
24 configurations) while one blind seed still finds the best configuration. With
three seeds this is a trend, not a firm result.

Three more seeds on the 4090 (six in all, each running all three arms) firm up
one conclusion and soften another. After 24 configurations the medians are
guided 0.065 s [0.064, 0.066], blind 0.077 s [0.064, 0.119] on the small space and
0.080 s [0.059, 0.119] on the large one; after 16, 0.086 s [0.086, 0.093], 0.099 s
[0.065, 0.469] and 0.142 s [0.078, 0.651]. The guided arm comes within 5\% of
0.064 s in 6 of 6 seeds, each blind arm in 2 of 6: the diagnosis makes the search
reliable. The larger space hurts the blind search clearly after 16
configurations and only slightly after 24, so the effect of space size is not
established. One blind seed on the large space beat every guided run: 0.059 s in
32 evaluations, from a finer cadence ($K=8$) and a cheaper step estimate (10 power
iterations instead of 30) on the working set. No rule proposed either change: the
cadence rule stopped at $K=32$ and no rule touched the number of power
iterations. Two rules were added for these, again a failure of the loop turned
into a rule, and the guided arm was rerun (six seeds each): with the rules as
first written (a condition error kept the cadence rule from ever firing), fixed,
and fixed with the guided arm limited to four proposals per expanded
configuration, rules for the dominant factor of $T=N\times t$ first. None moved the
median after 24 configurations (0.0660, 0.0648 and 0.0648 s against 0.0647 s), and
the 0.059 s configuration was reached at most once. A direct timing (ten
interleaved solves per configuration) shows why and settles whether the blind
result was noise. The guided best takes 0.0667 s [0.0633, 0.0718]; the blind best
0.0568 s [0.0528, 0.0686], a real gain; the guided best with only the cheaper step
estimate 0.0492 s [0.0468, 0.0534], the fastest configuration found so far
($3.4\times$ cuML's fastest run on this GPU); and the guided best with only the finer
cadence 0.0891 s [0.0821, 0.0978], worse. So the step-estimate rule is right and
points directly at the best configuration, but within 24 evaluations the guided
search spends its budget on the other proposals before it reaches it, while the
cadence rule is wrong and has been removed. Raising the budget to 32 configurations (six seeds,
step-estimate rule kept, cadence rule removed) leaves the median at 0.0649 s
[0.0629, 0.0662] and reaches the fastest configuration in no seed, and the records
show exactly why: on every seed the search first reaches the configuration it
must extend (the working set started at $p/16$ with $K=32$) at evaluation 21, and
expands it only in the last round, where the rules propose keeping momentum
(evaluation 30), a finer cadence (evaluation 31) and then the cheaper step
estimate, one evaluation past the budget. The search is deterministic apart from
timing noise, so the miss is structural: the rule is right, the order in which a
node's proposals are tried is not. Each new rule makes a fixed budget
thinner; ranking proposals by a calibrated prediction of their gain, the role of
CAKE's cost model, is the missing piece. Neither rule was used in the
generalization stage below, whose tuning procedure was fixed before they existed.

Working sets with dual-score selection are published practice (celer,
skglm, Blitz), so that part is an automatic rediscovery of known structure
exploitation, not a new algorithm; what is new is that the chain was derived
from measurements, per device, with the certificate unchanged. The ablation
below separates the ingredients, and the one carrying the speedup is not the
working set itself but the step taken from the Lipschitz constant of the
current working set, which that literature does not do: it is almost entirely
coordinate descent, where the step is per feature and computed once on the
whole matrix, so a restricted constant changes nothing there. Ours runs in
float64, cuML in float32. Compilation (2.7 to 3.9 s) depends only on $(n, p)$
and is excluded here, as is cuML's warm-up solve. That exclusion was misleading
when these rows were measured, because the programs were rebuilt on every call
and so the cost was paid on every solve; they are reused across
calls and across penalty weights, so a repeated solve and a path both pay it
once, and a first solve still pays it in full. cuML's time varied between runs of the same solve,
and the comparison uses its fastest. At looser tolerances cuML does not reach
the target (gap 0.32 at $10^{-2}$, $9\times10^{-3}$ at $10^{-3}$), so $10^{-4}$ is its
fastest passing setting. These rows are one instance at one penalty weight; the
ablation and the held-out stage below extend the comparison to six and ten.

\paragraph{Baseline precision.}
Our worker first ran cuML in float32 only (it cast the data before fitting). On
harder instances float32 cannot reach the target: on the epsilon test split at
$\alpha=0.9$, $\lambda=0.02\lambda_{\max}$ its gap stays at $1.4\times10^{-5}$ or
above for every tolerance from $10^{-4}$ to $10^{-8}$, and at $\alpha=0.5$,
$\lambda=0.02\lambda_{\max}$ at $1.4\times10^{-6}$ or above, while cuML in float64
certifies both (1.414 s, 317 iterations; 0.567 s, 127 iterations). On the
instance of the table above both precisions pass, float64 in 0.211 s in one
run, slower than float32's fastest (0.166 s), so the comparisons above stand. The
worker now takes the precision as an option, and every later comparison uses,
per instance, the faster of the two precisions that reaches the target. Counting
an instance where float32 fails as a win would reward our float64 arithmetic,
not our method.

\paragraph{The loop run end to end.}
Artifact: . Starting from
plain forward-backward with Nesterov momentum on the epsilon test split at
$\alpha=0.5$, $\lambda=0.1\lambda_{\max}$, which loses to cuML by $4.8\times$
($0.770$ s against $0.159$ s), the search reached $0.059$ s in four diagnosed
steps: screening ($0.590$ s), a working set ($0.144$ s), a smaller initial
working set ($0.081$ s), and a halved certificate cadence ($0.059$ s). That is
$2.67\times$ faster than the baseline it started behind. Each step is one
symptom mapped to one typed mutation, and the certificate is the full-problem
duality gap throughout, so nothing in the chain changes what is being certified.
The operator ingredients were proposed, admitted by the gate, run, and rejected
on measurement rather than by assumption: the anchored iteration could not
certify within twenty times the incumbent's iteration count and was stopped
there, and safeguarded Anderson was slower by a factor of fifty. On this cell
the gain is entirely on the screening axis.

Bounding the search was necessary and is worth recording, because two earlier
attempts at the same run exhausted their time budget. A step that loses badly is
followed once and then neither re-proposed from other parents nor expanded into
its own descendants, and a composed recipe is cut off once it exceeds twenty
times the best certified iteration count, which is a competitiveness bound
rather than a convergence claim. An earlier attempt also stopped one step short,
at $1.93\times$, for a reason that was our instrument and not the method: several
candidates in the last round failed for want of device memory, and those
failures were recorded as slow configurations rather than as failures to
measure. Distinguishing the two did not merely correct the record, it let the
search find the cadence step and a configuration $1.4\times$ faster than the one
it had settled for.

\subsection{Isolating the ingredients}
\label{app:log:ablation}

 The chain above mixes several changes, so
this stage varies one at a time. Ten instances of the epsilon test split,
$\alpha \in \{0.5, 0.9\}$ and $\lambda/\lambda_{\max} \in \{0.02, 0.05, 0.1, 0.2,
0.5\}$, five arms, seven repetitions each. The repetitions are interleaved, so
that repetition $k$ of every arm runs next to repetition $k$ of the others:
pairing arms by repetition index means nothing when each arm's repetitions run
in one block minutes apart, which is how an earlier version of this experiment
was arranged, and the paired intervals below are reported only because the
interleaved design earns them.

The arms are plain forward-backward with momentum; a working set whose step
comes from the Lipschitz constant of the whole matrix, with momentum either
reset or carried across changes of the set; and a working set whose step comes
from the Lipschitz constant of the current set, again with momentum reset or
carried. Three times are reported separately throughout, because they are
three different claims: the solve loop, a repeated solve (the loop plus
compilation, with compiled programs reused), and a first solve on an empty
cache.

\begin{center}
\small
\begin{tabular}{lcc}
\toprule
Contrast & geometric mean & $95\%$ interval \\
\midrule
reduced-set step against whole-matrix step, momentum reset & $2.058\times$ & $[1.635, 2.590]$ \\
reduced-set step against whole-matrix step, momentum carried & $1.840\times$ & $[1.426, 2.374]$ \\
momentum carried against reset, whole-matrix step & $1.171\times$ & $[1.103, 1.243]$ \\
momentum carried against reset, reduced-set step & $1.047\times$ & $[0.992, 1.105]$ \\
both against neither & $2.154\times$ & $[1.656, 2.801]$ \\
both against plain forward-backward & $11.687\times$ & $[8.227, 16.601]$ \\
\bottomrule
\end{tabular}
\end{center}

The step taken from the Lipschitz constant of the current working set is what
carries the result: its interval excludes one on every instance and across
them. Carrying momentum is real only while the step comes from the whole
matrix; once the reduced-set step is in place the interval straddles one and
the within-instance comparison is inconclusive on seven of the ten. The reason
is visible in the counts: the reduced-set step cuts a solve from 288
evaluations and three re-plans to 96 and two, and momentum can only be worth
what the re-plans would otherwise destroy. Against cuML, taking per instance
the faster of its two precisions that reaches the target and its fastest run of
that precision, the best arm is $2.28\times$ with a $95\%$ interval of $[1.73,
3.01]$ and faster on ten of ten for the solve loop and for a repeated solve,
and $0.06\times$ and faster on none for a first solve. Running every arm in
float32, which matches the precision cuML wins with on most of these
instances, gives $2.43\times$ and does not change any of the contrasts.

One negative is worth stating because it was the hypothesis this grid was run
to test. The advantage over cuML grows with the penalty weight, so it would be
convenient if the ingredient that carries the speedup grew with it too. It does
not. At $\alpha = 0.5$ the gain from the reduced-set step falls from
$1.96\times$ at $0.02\lambda_{\max}$ to $1.13\times$ at $0.5\lambda_{\max}$,
while the advantage over the baseline rises across that same range, so the two
move in opposite directions; at $\alpha = 0.9$ the gain is not monotone at all
and peaks in the middle. The two effects have different explanations. Heavily
penalized instances are already short, thirty-two evaluations at
$0.5\lambda_{\max}$, so a larger step has little left to save and the advantage
there comes from the working set being small, which is a cost-per-iteration
effect; lightly penalized instances give the larger step the most room in
relative terms while remaining hard in absolute terms, which is where the
baseline closes. The ingredient acts on the iteration count and the regime
boundary is set by the cost per iteration.

\subsection{Generalization: a frozen portfolio on held-out instances}
\label{app:log:general}

 Following
the separation CAKE (arXiv 2608.12629) makes between tuning one shape and building a
library, the domain and its split were written to before any
tuning: four real LIBSVM datasets (epsilon training split,
$4\times10^5\times2000$; epsilon test split, $10^5\times2000$; gisette, $6000\times5000$;
covtype, $581012\times54$), $\lambda/\lambda_{\max}\in\{0.5, 0.2, 0.1, 0.05, 0.02\}$ and
$\alpha\in\{0.5, 0.9\}$. The 18 tuning instances are epsilon training, gisette and
covtype at 0.5, 0.1 and 0.02; the 22 held-out instances are the same datasets at 0.2
and 0.05 (unseen regularization) and every epsilon test instance (an unseen shard of a
tuned source). The guided loop ran on every tuning instance on the Blackwell, step
size timed, including the 9 on which cuML in float32 never reaches the target (there
it is tuned against cuML's closest run, which only feeds the rules). Every distinct
best configuration was then timed on every tuning instance, and a portfolio of at most
three recipes with a dispatch tree of depth two on declared features ($\log n$,
$\log p$, $\lambda/\lambda_{\max}$, $\alpha$) was fitted to the geometric-mean time on the
tuning instances only, then frozen with a content hash (\texttt{ae2a840e\ldots}) before
any held-out instance ran.

The fitted portfolio has two recipes and one split, on $\lambda/\lambda_{\max}\le0.06$:
below it, a working set started at $p/64$ with screening, buckets and $K=64$; above it,
the same with $K=16$ and float32 storage. On the tuning instances its geometric-mean
time is 0.071 s, against 0.084 s for the single best recipe and 0.060 s for an oracle
that picks the best candidate per instance.

Held-out result against cuML in both precisions (three runs per instance, fastest
shown; our time is the median of three solves after an untimed warm-up; the last
column divides the faster cuML precision that reaches the target by our time):

\begin{center}
\small
\begin{tabular}{lllrrrr}
\toprule
dataset & $\alpha$ & $\lambda/\lambda_{\max}$ & ours (s) & cuML float32 (s) & cuML float64 (s) & speedup \\
\midrule
covtype & 0.5 & 0.05 & 0.1102 & fails ($\ge4\times10^{-6}$) & 0.3258 & 2.96 \\
covtype & 0.5 & 0.2 & 0.0552 & 0.0698 & 0.0778 & 1.26 \\
covtype & 0.9 & 0.05 & 0.2265 & fails ($\ge3\times10^{-5}$) & 0.6153 & 2.72 \\
covtype & 0.9 & 0.2 & 0.0670 & fails ($\ge2\times10^{-5}$) & 0.1330 & 1.98 \\
epsilon test & 0.5 & 0.02 & 0.2380 & 0.3296 & 0.3679 & 1.39 \\
epsilon test & 0.5 & 0.05 & 0.1806 & 0.2085 & 0.2931 & 1.15 \\
epsilon test & 0.5 & 0.1 & 0.0499 & 0.1296 & 0.1587 & 2.60 \\
epsilon test & 0.5 & 0.2 & 0.0388 & 0.1372 & 0.0799 & 2.06 \\
epsilon test & 0.5 & 0.5 & 0.0103 & 0.1465 & 0.0638 & 6.18 \\
epsilon test & 0.9 & 0.02 & 0.2730 & fails ($\ge5\times10^{-6}$) & 0.8840 & 3.24 \\
epsilon test & 0.9 & 0.05 & 0.1286 & fails ($\ge1\times10^{-6}$) & 0.4922 & 3.83 \\
epsilon test & 0.9 & 0.1 & 0.0697 & 0.3237 & 0.3375 & 4.65 \\
epsilon test & 0.9 & 0.2 & 0.0327 & 0.1879 & 0.2496 & 5.74 \\
epsilon test & 0.9 & 0.5 & 0.0130 & 0.1180 & 0.0786 & 6.03 \\
epsilon train & 0.5 & 0.05 & 0.6301 & 0.6861 & 0.8139 & 1.09 \\
epsilon train & 0.5 & 0.2 & 0.1877 & 0.4417 & 0.1857 & 0.99 \\
epsilon train & 0.9 & 0.05 & 0.4403 & fails ($\ge3\times10^{-6}$) & 1.7208 & 3.91 \\
epsilon train & 0.9 & 0.2 & 0.1326 & 0.6295 & 0.6359 & 4.75 \\
gisette & 0.5 & 0.05 & 0.0313 & 0.1160 & 0.2050 & 3.70 \\
gisette & 0.5 & 0.2 & 0.0194 & 0.0585 & 0.1072 & 3.02 \\
gisette & 0.9 & 0.05 & 0.0443 & fails ($\ge3\times10^{-6}$) & 0.3383 & 7.64 \\
gisette & 0.9 & 0.2 & 0.0178 & 0.0800 & 0.1352 & 4.49 \\
\bottomrule
\end{tabular}
\end{center}

Our dispatched recipe certifies all 22 held-out instances, and cuML in float64
reaches the target on all 22 as well, so every instance has a baseline. Against
the faster cuML precision that reaches the target on each instance, the frozen
portfolio is faster on 21 of 22, by a geometric mean of $2.91\times$ against cuML's
fastest run and $3.13\times$ against its median (range $0.99\times$ to $7.64\times$;
the one loss, epsilon training at $\alpha=0.5$, $\lambda=0.2\lambda_{\max}$, is 0.1877 s
against 0.1857 s). Against float32 alone, on the 15 instances where it reaches the
target, the figures are $3.21\times$ and $3.82\times$ (15 of 15); against float64 alone,
$3.33\times$ and $3.42\times$ (21 of 22). For reference, the single best recipe gives
$2.37\times$ (20 of 22) and plain FISTA $0.36\times$ (3 of 22) against the faster cuML
precision, so the gain comes from the recipes the loop found and from dispatching
between two of them, not from FISTA. The tuning used the rules as they stood before
the two rules of Section~\ref{app:log:diagnose} were added; the cheaper step estimate
those runs identified is not in this portfolio.

This comparison covers one regularization family (elastic-net logistic), four dense datasets and
one GPU type; cuML's times vary between runs (the table shows the fastest of three);
ours runs in float64 or float32 storage with a float64
certificate, and cuML was given both precisions.

The exclusion of compilation above is not the symmetric statement it appears to be,
and the numbers in this subsection must be read with the following correction. Our
compilation is about three seconds against solves of $0.01$ to $0.4$ s, and it was re-paid on every call, because the compiled programs were held in a
dictionary local to the solve and discarded when it returned. The excluded cuML
warm-up is a two-iteration fit on a small slice. So the figure reported above was the
solve loop alone, and no caller could obtain it: every solve, including every point of
a regularization path, paid the compilation again. Programs are now reused across
calls, keyed on a fingerprint of the labels as well as the shapes, since the labels and
the two penalty parameters are constants in the compiled program. A repeated solve at
one penalty weight now compiles once, and so does a path over several weights: the two
penalty parameters were made arguments of the compiled program rather than constants
in it. Solving a five-point path now costs one compilation instead of five, and the
solutions are unchanged to the last bit.

Three times are therefore distinct and they are three different claims: a first solve
on an empty cache, a repeated solve, and the solve loop. Re-measured under that
accounting, on the ten held-out instances that can still be measured, the dispatched
portfolio is $2.39\times$ the better cuML precision's fastest run, with a $95\%$
interval over the ten instances of $1.40$ to $4.07$, and faster on eight of ten; the
repeated solve is identical, and these records do not measure a first solve, because
the warm-up absorbs the compilation. Variation between runs is small next to variation
between instances, which is what makes an interval over instances the right summary
here: our three repetitions span at most $1.06\times$ and cuML's at most $1.57\times$,
and the comparison takes cuML's fastest run rather than its median, which favours the
baseline. It is an interval over instances and not a paired test between the two
solvers, whose repetitions are independent runs with no matched pairing. The advantage grows monotonically with
the penalty weight, from $1.50\times$ at $\lambda = 0.02\lambda_{\max}$ to
$7.00\times$ at $0.5\lambda_{\max}$, and both losses fall at the weakest penalties,
where the solution is densest and screening removes least. The eight-fold spread is
visible in the iteration counts, which run from 16 at $0.5\lambda_{\max}$ to 448 at
$0.02\lambda_{\max}$. The claim this supports is therefore about well-regularized
instances of the family, with a crossover near $0.05\lambda_{\max}$ at $\alpha=0.5$,
not about the family as a whole.

Only ten of the twenty-two held-out instances could be re-measured. Two of the four
datasets are no longer on disk and are absent from the data manifest, and the largest
does not fit in this GPU's memory beside the language-model server, needing
$5.98$ GiB for the design matrix alone. The earlier $2.91\times$ over twenty-two
instances and the $2.39\times$ here are not the same quantity measured twice.

\subsection{Variable metric ideas considered for QP}
\label{app:log:varmetric}

Source: J.F.~Bonnans, J.Ch.~Gilbert, C.~Lemar\'echal and C.A.~Sagastiz\'abal, A
family of variable metric proximal methods, Mathematical Programming 68
(1995) 15--47, read in full. The first idea is now implemented and measured
(Section~\ref{app:log:adaptive}). The quasi-Newton primal metric is implemented as
the builder \texttt{condat\_vu\_metric} and measured on all 137 QP instances: it does not
help. With ten Krylov secant pairs it solves 55 instances like the scalar step (one
gained, one lost) and needs $1.20\times$ the iterations and $1.52\times$ the time where
both solve; built from warm-up iterates it solves 53 (two gained, four lost), with
$1.16\times$ the iterations. The remaining ideas are recorded with the gate rule they
would need.

\begin{itemize}
\item A gate for adapting the metric. The paper's global result allows any
 sequence of positive definite metrics $M_n$ in the proximal step, provided a
 sufficient-decrease test holds and $\sum_n \lambda_{\min}(M_n^{-1}) = \infty$
 (for BFGS metrics this follows from $\operatorname{tr} M_n \le (n+1)C$). For
 the splitting schemes in our genome the matching condition is the variable
 metric quasi-Fej\'er rule of Combettes and V\~u (2013, 2014): uniform bounds
 $\alpha I \preceq M_n \preceq \beta I$ and $M_{n+1} \preceq (1+\eta_n) M_n$
 with $\sum_n \eta_n < \infty$. This is the missing certificate for the
 adaptive step ratio: PDLP's primal weight and OSQP's adaptive $\rho$ are
 scalar varying metrics adapted without such a budget.
\item A quasi-Newton primal metric. With secant pairs $s_n = x_{n+1}-x_n$ and
 $y_n = \nabla f(x_{n+1}) - \nabla f(x_n)$, which for a QP equal $P s_n$
 exactly, a limited-memory BFGS metric can drive the primal step. The
 structural fact that makes this cheap: in the OSQP form of a QP the
 separable primal term is zero, so the primal update
 $x^+ = x - M_n^{-1}(\nabla f(x) + A^\top u)$ needs no proximal map in the
 non-diagonal metric, only an application of $M_n^{-1}$. The step condition
 in the new metric must be re-checked (power iteration on $M_n^{-1}P$ and
 $A M_n^{-1} A^\top$), which is affordable only at chunk boundaries: metric
 updates would happen at the certificate cadence $K$, the one gene already
 known to couple $N$ and $t$.
\item Reflection with a matched metric. When the metric equals the Hessian of
 a quadratic, the proximal step moves only half way to the minimizer, and the
 step $x_{n+1} = x_n + 2(x^p_n - x_n)$ recovers it; the paper's line search
 prefers this step and falls back when a descent test fails. In our genome
 this is relaxation $\rho=2$ combined with the safeguarded switch wrapper.
 Reflected Halpern restarts without a matched metric did not help on
 transport (Section~\ref{app:log:ot}).
\item A relative acceptance test for inexact proximal steps. An approximate
 proximal point is accepted when a lower model of the objective is accurate
 enough there, a relative test that needs no summable error sequence. This
 is the stopping rule the conjugate-gradient option of a factorization-based
 proximal step would need, and it connects to the relative-error criterion
 of Solodov and Svaiter in the knowledge base. It is not new for QP:
 rPDHCG solves the primal subproblems of PDHG inexactly by conjugate
 gradients, and a 2026 GPU conic-QP solver uses a relative proximal-error
 budget for restarted PDHG (arXiv 2608.09159).
\item Prior use, from one search each: a quasi-Newton primal metric appears
 in general proximal splitting (arXiv 1801.08691, 2209.14019) but was not
 found in LP or QP PDHG practice, so it is a frontier candidate pending a
 wider search; reflection is already used (reflected Halpern PDHG for LP,
 arXiv 2407.16144).
\end{itemize}
